\section{Pair of genealogical readers over the complete atlas}
\label{sec:operadores-masa-two-way-rev11}
\index{mass!bidirectional operators}
\index{atlas!chronological and dodecaphasic readers}

The enriched route possesses two complementary internal representations,
\[
 \Gamma\longmapsto
 \bigl(\gamma_{108}(\Gamma),\tau_{12}(\Gamma)\bigr)
 \in\mathbb F_3^{108}\times\{-1,0,+1\}^{12}.
\]
The first preserves the complete tritic history; the second preserves the
signed dodecaphase record. Each projection induces a natural reader, and both
are evaluated before opening masses, PDG labels, or residuals.

\HMTClaim{MD-R-MASS-TWOWAY-READER-PAIR}
\subsection{Pair of genealogical readers \(\mathbf K(\Gamma)=(K_D(\Gamma),K_\Omega(\Gamma))\): domain, codomain, and scalar evaluations}

Using the notation of
\cref{eq:KD-rev11,eq:kappaD-rev11,eq:kappa-omega-rev11}, we define
\begin{equation}
 \mathbf K(\Gamma)
 :=\bigl(K_D(\Gamma),K_\Omega(\Gamma)\bigr)
 \in\left(\mathbb Z^2\times\tfrac12\mathbb Z\right)
 \times\mathbb Z^3,
 \label{eq:lector-conjunto-masas-rev11}
\end{equation}
where
\[
 K_D=\left(D_{27}+D_{36},D_1,\frac{D_4-D_3}{2}\right),
 \qquad
 K_\Omega=(\Omega_{90\mid120},54,P_6).
\]
Scalar evaluations are
\begin{align}
 \kappa_D(\Gamma)
 &=D_{27}+D_{36}-\alpha_{\rm HMT}D_1
 +\frac{\Delta_4}{2}(D_4-D_3),
 \label{eq:kappaD-general-rev11}\\
\kappa_\Omega(\Gamma)
&=\Omega_{90\mid120}(\tau_\Gamma)
-54\alpha_{\rm HMT}+P_6(\tau_\Gamma)\Delta_4.
\label{eq:kappaOmega-general-rev11}
\end{align}
The second coordinate requires a causal distinction. In APP, the primordial spectral jump is \(9-5=4\); the local coefficient \(2\), the compression \(6=51-45\), and the two-dimensional unfolding \(54=9\cdot6=459-405\) are later typed invariants, in accordance with the \cref{thm:app-jerarquia-cuatro-dos-seis-cincuentaycuatro}. By contrast, in the chronological displacement reader, the electron coordinate is \(B_D=D_1(\Gamma_e)=54\): it counts immediate boundaries and coincides with the kinematic periodicity of order \(54\) of the CT108 endomorphism. The fixed \(54\) of \(K_\Omega\) is the common coordinate transferred when both readers are made compatible in the electron reduction. Its numerical equality with the global APP defect constitutes a subsequent cross-check; it does not make the number \(54\) the primordial difference between addition and multiplication.

The equality \(4\cdot90=3\cdot120=360\) explains the common genealogy of
the two readers. \(K_D\) compares displacements of the tritic word
\(\gamma_{108}\);
\(K_\Omega\) compares the energies of windows in the twelve-phase word. The
coincidence of origin does not identify their information quotients.
In the certified domain \(\Gamma_{324}^{\rm or}\), half-turn periodicity
makes all the \(D_k\) even, so that
\((D_4-D_3)/2\in\mathbb Z\) and \(K_D\in\mathbb Z^3\). Outside that domain,
the third component retains type \(\tfrac12\mathbb Z\).

\subsection{Naturalness of bidirectional operators and type preservation}

The reader \(D_k\) is invariant under time translation, reversal, and every
bijective permutation of the tritic alphabet. On the declared routes,
simultaneously reversing both orientations acts by a translation of
27 discrete steps and preserves \(K_D\). The reader \(K_\Omega\) is invariant under
dodecaphase rotation and reversal and under
\(C_M(\tau)=-\operatorname{rev}(\tau)\). Its quadratic form
\[
 \Omega(\tau)=\langle\tau,
 (4W_3^*W_3-3W_4^*W_4)\tau\rangle
\]
possesses integer bilinear polarization and satisfies the 2-cocycle identity.

The global counters of turn, sheet change, phase repetition, prefix change, and return are constant over the 324 routes. Consequently, every functional that factors exclusively through those six counters is constant and cancels in ratios. The readers in \eqref{eq:kappaD-general-rev11} and \eqref{eq:kappaOmega-general-rev11} act on the local history and pass that nonfactorization control.

\begin{ablation}[Necessity of the two sheets in the electronic route]
\label{abl:hojas-lector-two-way-rev11}
The chronological reader produces
\[
 (A_D,B_D,C_D)=
 \begin{cases}
 (80,54,6),&\substack{\text{joint chronological word of both sheets;}\\
 B_D=D_1(\Gamma_e)=54},\\
 (80,84,-8),&\text{additive projection},\\
 (80,60,13),&\text{multiplicative projection}.
 \end{cases}
\]
The complete electronic triple requires the joint history of the two arithmetic sheets.
\end{ablation}

\subsection{Exact Censuses over the 324 Routes}

The chronological reader produces fourteen profiles. The following table records their
complete census; \(A_D=D_{27}+D_{36}\), \(B_D=D_1\), and
\(C_D=(D_4-D_3)/2\).

\begin{longtable}{@{}rrr r@{}}
\caption{Exact census of the fourteen profiles of the chronological reader over the
324 oriented routes.}
\label{tab:censo-perfiles-kd-rev11}\\
\toprule
\(A_D\)&\(B_D\)&\(C_D\)&multiplicity\\
\midrule
\endfirsthead
\toprule
\(A_D\)&\(B_D\)&\(C_D\)&multiplicity\\
\midrule
\endhead
0&24&0&18\\
40&24&2&36\\
40&54&6&18\\
40&78&27&36\\
40&84&30&36\\
60&78&25&18\\
60&84&28&18\\
80&54&6&18\\
80&54&7&18\\
80&66&2&18\\
80&66&3&36\\
80&66&4&18\\
100&66&0&18\\
100&66&2&18\\
\midrule
\multicolumn{3}{r}{total}&324\\
\bottomrule
\end{longtable}

The dodecaphasic reader produces nine profiles:

\begin{table}[H]
\centering
\caption{Exact census of the nine profiles of the dodecaphasic reader over
the 324 oriented routes.}
\label{tab:censo-perfiles-komega-rev11}
\begin{tabular}{@{}rrr r@{}}
\toprule
\(\Omega\)&\(B_\Omega\)&\(P_6\)&multiplicity\\
\midrule
80&54&6&198\\
56&54&5&68\\
48&54&5&34\\
36&54&4&8\\
20&54&4&4\\
4&54&2&4\\
40&54&4&4\\
24&54&4&2\\
24&54&2&2\\
\midrule
\multicolumn{3}{r}{total}&324\\
\bottomrule
\end{tabular}
\end{table}

The readers coincide as triples on (18) routes, always with the value \((80,54,6)\), and the pair \(\mathbf K\) assumes (40) values. In the canonical section \(B1\) of 81 positions, \(K_D\) is constant in 41 of the 56 families and refines the quotient into 77 classes; \(K_\Omega\) is constant in 50 and produces 62 classes; the pair produces 78. In the lift to 324 orientations, \(K_D\) produces 146 classes and is nonconstant in every family; \(K_\Omega\) produces 100 and descends in 17 families; the pair produces 151 and is nonconstant in all of them. These censuses prove that the two projections retain distinct memories and determine the domain of each descent claim.

\subsection{Vector traces and joint spectra of the \(13\) multisections}

In the canonical section \(B1\), we denote by
\(\ell_j,\nu_j,d_j,u_j\) the generation-\(j\) multisections of the
charged-lepton, neutrino, down-type-quark, and up-type-quark
sectors, respectively. For each reader
\(r\in\{D,\Omega\}\) and each component \(a\in\{1,2,3\}\), we define the
diagonal operator
\[
 K_{M,a}^r e_\Gamma=K_{r,a}(\Gamma)e_\Gamma,
 \qquad
 \mathbf K_M^r=(K_{M,1}^r,K_{M,2}^r,K_{M,3}^r).
\]
The three operators of \(\mathbf K_M^r\) are self-adjoint and commute.
Therefore, their normalized trace is understood componentwise,
\[
 \overline{\mathbf K}_M^r
 =\frac{1}{\dim M}
 \bigl(\operatorname{Tr}K_{M,1}^r,
       \operatorname{Tr}K_{M,2}^r,
       \operatorname{Tr}K_{M,3}^r\bigr),
\]
and its spectrum is the joint spectrum, that is, the distribution with
multiplicity of the simultaneous triples. The thirteen vector traces are:

\begin{longtable}{@{}c r c c@{}}
\caption{Normalized vector traces of the readers \(K_D\) and
\(K_\Omega\) on the thirteen canonical multisections.}
\label{tab:trazas-multisecciones-two-way-rev11}\\
\toprule
\(M\)&\(\dim M\)&\(\overline{\mathbf K}_M^D\)&\(\overline{\mathbf K}_M^\Omega\)\\
\midrule
\endfirsthead
\toprule
\(M\)&\(\dim M\)&\(\overline{\mathbf K}_M^D\)&\(\overline{\mathbf K}_M^\Omega\)\\
\midrule
\endhead
\(\ell_0\)&1&\((80,54,6)\)&\((80,54,6)\)\\
\(\ell_2\)&8&\((65,249/4,17/2)\)&\((80,54,6)\)\\
\(\ell_4\)&16&\((225/4,489/8,21/2)\)&\((141/2,54,11/2)\)\\
\(\nu_1\)&2&\((40,51,29/2)\)&\((60,54,5)\)\\
\(\nu_2\)&4&\((45,63,29/2)\)&\((61,54,5)\)\\
\(\nu_3\)&6&\((170/3,55,34/3)\)&\((206/3,54,11/2)\)\\
\(\nu_4\)&8&\((105/2,237/4,95/8)\)&\((143/2,54,45/8)\)\\
\(d_1\)&2&\((40,51,29/2)\)&\((60,54,5)\)\\
\(d_2\)&4&\((45,63,57/4)\)&\((61,54,5)\)\\
\(d_3\)&6&\((170/3,55,23/2)\)&\((218/3,54,17/3)\)\\
\(d_4\)&8&\((105/2,237/4,49/4)\)&\((149/2,54,23/4)\)\\
\(u_1\)&4&\((65,66,9)\)&\((80,54,6)\)\\
\(u_3\)&12&\((205/3,119/2,113/12)\)&\((74,54,23/4)\)\\
\bottomrule
\end{longtable}

The notation \((a,b,c)^{\times n}\) indicates multiplicity \(n\). The complete
joint spectrum of \(\mathbf K_D\) is:

\begingroup\footnotesize
\begin{longtable}{@{}c p{.82\textwidth}@{}}
\caption{Complete joint spectra of the reader \(K_D\) in the thirteen
canonical multisections.}
\\
\toprule
\(M\)&\(\operatorname{Spec}_{\rm conj}\mathbf K_M^D\)\\
\midrule
\endhead
\(\ell_0\)&\((80,54,6)^{\times1}\)\\
\(\ell_2\)&\((0,24,0)^{\times1};(40,78,27)^{\times2};(80,54,7)^{\times1};(80,66,2)^{\times1};(80,66,3)^{\times1};(100,66,0)^{\times1};(100,66,2)^{\times1}\)\\
\(\ell_4\)&\((0,24,0)^{\times1}\); \((40,24,2)^{\times2}\); \((40,54,6)^{\times2}\); \((40,84,30)^{\times2}\); \((60,78,25)^{\times3}\); \((80,66,2)^{\times2}\); \((80,66,3)^{\times3}\); \((80,66,4)^{\times1}\)\\
\(\nu_1\)&\((40,24,2)^{\times1};(40,78,27)^{\times1}\)\\
\(\nu_2\)&\((0,24,0)^{\times1};(40,84,30)^{\times1};(60,78,25)^{\times1};(80,66,3)^{\times1}\)\\
\(\nu_3\)&\((40,24,2)^{\times2};(40,78,27)^{\times1};(60,84,28)^{\times1};(80,54,6)^{\times1};(80,66,3)^{\times1}\)\\
\(\nu_4\)&\((0,24,0)^{\times1};(40,24,2)^{\times1};(40,78,27)^{\times2};(40,84,30)^{\times1};(80,54,7)^{\times1};(80,66,2)^{\times1};(100,66,0)^{\times1}\)\\
\(d_1\)&\((40,24,2)^{\times1};(40,78,27)^{\times1}\)\\
\(d_2\)&\((0,24,0)^{\times1};(40,84,30)^{\times1};(60,78,25)^{\times1};(80,66,2)^{\times1}\)\\
\(d_3\)&\((40,24,2)^{\times2};(40,78,27)^{\times1};(60,84,28)^{\times1};(80,54,6)^{\times1};(80,66,4)^{\times1}\)\\
\(d_4\)&\((0,24,0)^{\times1};(40,24,2)^{\times1};(40,78,27)^{\times2};(40,84,30)^{\times1};(80,54,7)^{\times1};(80,66,3)^{\times1};(100,66,2)^{\times1}\)\\
\(u_1\)&\((40,54,6)^{\times1};(60,78,25)^{\times1};(80,66,2)^{\times1};(80,66,3)^{\times1}\)\\
\(u_3\)&\((40,24,2)^{\times2}\); \((40,78,27)^{\times2}\); \((60,84,28)^{\times1}\); \((80,54,6)^{\times3}\); \((80,66,3)^{\times1}\); \((80,66,4)^{\times1}\); \((100,66,0)^{\times1}\); \((100,66,2)^{\times1}\)\\
\bottomrule
\end{longtable}
\label{tab:espectros-kd-multisecciones-rev11}
\endgroup

The complete joint spectrum of \(\mathbf K_\Omega\) on the same fibres is:

\begingroup\footnotesize
\begin{longtable}{@{}c p{.82\textwidth}@{}}
\caption{Complete joint spectra of the reader \(K_\Omega\) in the thirteen
canonical multisections.}
\label{tab:espectros-komega-multisecciones-rev11}\\
\toprule
\(M\)&\(\operatorname{Spec}_{\rm conj}\mathbf K_M^\Omega\)\\
\midrule
\endfirsthead
\toprule
\(M\)&\(\operatorname{Spec}_{\rm conj}\mathbf K_M^\Omega\)\\
\midrule
\endhead
\(\ell_0\)&\((80,54,6)^{\times1}\)\\
\(\ell_2\)&\((80,54,6)^{\times8}\)\\
\(\ell_4\)&\((24,54,2)^{\times1};(56,54,5)^{\times4};(80,54,6)^{\times11}\)\\
\(\nu_1\)&\((40,54,4)^{\times1};(80,54,6)^{\times1}\)\\
\(\nu_2\)&\((4,54,2)^{\times1};(80,54,6)^{\times3}\)\\
\(\nu_3\)&\((36,54,4)^{\times1};(56,54,5)^{\times1};(80,54,6)^{\times4}\)\\
\(\nu_4\)&\((36,54,4)^{\times1};(56,54,5)^{\times1};(80,54,6)^{\times6}\)\\
\(d_1\)&\((40,54,4)^{\times1};(80,54,6)^{\times1}\)\\
\(d_2\)&\((4,54,2)^{\times1};(80,54,6)^{\times3}\)\\
\(d_3\)&\((36,54,4)^{\times1};(80,54,6)^{\times5}\)\\
\(d_4\)&\((36,54,4)^{\times1};(80,54,6)^{\times7}\)\\
\(u_1\)&\((80,54,6)^{\times4}\)\\
\(u_3\)&\((56,54,5)^{\times3};(80,54,6)^{\times9}\)\\
\bottomrule
\end{longtable}
\endgroup

The traces may coincide between multisections of different type; the type,
multiplicity, and full spectrum remain data of the operator.
The certificate reproduces these tables row by row.

\HMTClaim{MD-R-MASS-FIBRE-OPERATORS}
\subsection{Fiber Operators and Naturality of Base Change}

Distinguimos the fiber canonical \(F_M^{B1}\), formed by the positions of
the multisection \(M\) with its orientation published, and the fiber oriented
\(F_M^{\rm or}\), formed by their four lifts by position. For
\(s\in\{B1,{\rm or}\}\), let
\(\mathcal H_{M,s}=\bigoplus_{\Gamma\in F_M^s}\mathbb C e_\Gamma\). We define
\begin{equation}
 B_{M,s}^r e_\Gamma=\kappa_r(\Gamma)e_\Gamma,
 \qquad
 R_{\rm act}^{B_{M,s}^r}
 =\exp\bigl[(\log R_{\rm act})B_{M,s}^r\bigr],
 \qquad r\in\{D,\Omega\}.
 \label{eq:operadores-fibra-two-way-rev11}
\end{equation}
Each \(B_{M,s}^r\) is self-adjoint and diagonal in the route basis, so
\(R_{\rm act}^{B_{M,s}^r}\) is positive. For a positive basal mass
\(\widehat M_{M,s}^{(0)}\), the operatorial form conditioned on each reader is
\begin{equation}
 \widehat M_{M,s}^{\beta,r}
 =R_{\rm act}^{B_{M,s}^r/2}\,
  \widehat M_{M,s}^{(0)}\,
  R_{\rm act}^{B_{M,s}^r/2}.
 \label{eq:ley-operatoria-fibra-rev11}
\end{equation}
This expression is positive without imposing commutativity. When
\([\widehat M_{M,s}^{(0)},B_{M,s}^r]=0\), it reduces to
\(\widehat M_{M,s}^{(0)}R_{\rm act}^{B_{M,s}^r}\).
A unitary change of basis \(U\) transforms
\(B_{M,s}^r\mapsto UB_{M,s}^rU^*\) and preserves spectrum, trace, determinant, and
characteristic polynomial. The canonical electronic fiber \(F_e^{B1}\) has
rank one and reduces to the scalar of the
\cref{thm:reduccion-electron-two-way-rev11}.
Its oriented lift has rank four: the routes \((++)\) and \((--)\)
share the electronic profile, whereas the mixed orientations preserve
distinct profiles.

The normalized determinant
\[
 \operatorname{ndet}(R_{\rm act}^{B_{M,s}^r})
 =R_{\rm act}^{\operatorname{Tr}(B_{M,s}^r)/\dim F_M^s}
\]
is natural under permutation of the fiber. Logarithmic additivity distinguishes it as a geometric mean, but it is not used as a substitute for the physical realization. The exact refinement output remains the pair \((B_{M,s}^D,B_{M,s}^\Omega)\); the scalar section of a species is obtained through the flavor--generation--route morphism already constructed in \cref{eq:morfismo-sabor-ruta-cerrado-rev2,eq:realizacion-fisica-sabor-ruta-cerrada-rev2}: it selects the persistent route through non-metrological data and then applies the character of its signature. An arbitrary scalar functional \(S(B_{M,s}^D,B_{M,s}^\Omega)\) does not define that selection and remains only as a check that the target mass is not reintroduced.

\subsection{Absolute scale, relative correction and sensitivities}

The genealogy of the action line is written
\[
 \Sigma_H=\varphi\alpha_{\rm HMT}^{16},
 \qquad \operatorname{ord}_{10}(\Sigma_H)=-34,
\]
\[
 \delta_{2w}=H_5-\eta_{\rm ret}
 =\frac{90}{\pi}\mathscr C_\pi(\alpha_{\rm HMT}),
 \qquad R_{\rm act}=\frac{H_5}{\eta_{\rm ret}},
\]
\[
 \hbar_{\rm pre}=H_5\,10^{-34}\mathcal U_S,
 \qquad
 \hbar_{\rm ret}=\eta_{\rm ret}\,10^{-34}\mathcal U_S.
\]
The factor \(10^{-34}\mathcal U_S\) fixes the absolute section and cancels in
ratios. For the dodecaphase reader,
\begin{equation}
 \log\!\left[
 \frac{(m_X^{\beta,\Omega}/m_e^{\beta})}
 {(m_X^{(0)}/m_e^{(0)})}\right]
 =\bigl[(\Omega_X-80)+(P_{6,X}-6)\Delta_4\bigr]
 \log R_{\rm act}.
 \label{eq:razon-masa-omega-rev11}
\end{equation}
The common term \(-54\alpha_{\rm HMT}\) also cancels. For the chronological
reader,
\begin{equation}
 \log\!\left[
 \frac{(m_X^{\beta,D}/m_e^{\beta})}
 {(m_X^{(0)}/m_e^{(0)})}\right]
 =\bigl[\kappa_D(\Gamma_X)-\kappa_e\bigr]\log R_{\rm act}.
 \label{eq:razon-masa-D-rev11}
\end{equation}
With \(\log R_{\rm act}=6.932817628081925\ldots\times10^{-10}\),
\[
 \frac{\partial\log m}{\partial\Omega}=6.9328\times10^{-10},
 \quad
 \frac{\partial\log m}{\partial P_6}
 =\Delta_4\log R_{\rm act}=7.7090\times10^{-14},
\]
\[
 \frac{\partial\log m}{\partial\alpha}
 =-54\log R_{\rm act}=-3.74372\times10^{-8},
 \quad
 \left.\frac{\partial\log m}{\partial\Delta_4}\right|_e
 =6\log R_{\rm act}=4.15969\times10^{-9}.
\]
These derivatives are model sensitivities and remain separate from the
experimental uncertainty of a contrast mass.

\subsection{Physical pairing and null sector}

The typed pairing follows the chain
\[
 \begin{aligned}
 &\text{typed physical record}
 \longrightarrow\text{multisection or representation}
 \longrightarrow\text{fiber of routes}\longrightarrow\text{record}\\
 &\longrightarrow\text{signature and character}
 \longrightarrow(B_{M,s}^D,B_{M,s}^\Omega)
 \longrightarrow\text{dimensional chart}
 \longrightarrow\text{external comparison}.
 \end{aligned}
\]
The typed physical inventory distinguishes registers by object, representation, and codomain; the 324 routes are the internal engine; the 471 rows of the PDG catalogue are external properties. These cardinalities belong to different domains. The arrow is constructed from non-metrological physical descriptors, the multisection, and the route; a search by proximity to a mass belongs to the calibration regime.

The null sector bifurcates before the positive character. For photon and gluons,
\[
 \kappa_{\rm null}=\mathrm{N/A},
 \qquad m_{\rm null}=0.
\]
The expression \(R_{\rm act}^{0}=1\) is the unit of the positive character. Composite objects
concatenate records before evaluating the signature; topological sectors preserve fusion and
braiding; a virtual section is classified by its finite horizon of
prolongation.

\subsection{Independence of the path constructor with respect to masses and closure of the operator}

The structural verification uses exclusively the route census, the
register \(81\times108\), and the transducer of the 81 positions. It verifies their
row-by-row coupling, the coverage
\(81/56/13/324\), the naturality relations, the electronic triple \((80,54,6)\), the
censuses (14) and (9), the match (18/324), the
(40) pairs, differentiated descents in \(B1\) and in the 324 orientations,
periodicity 54, inversion by shift 27, and sheet ablations. The
null sector belongs to the domain exterior to the positive character. The calculation is
performed before introducing masses, CODATA, PDG, metrological residuals, or
nominal labels.

This certificate does not evaluate \(\alpha_{\rm HMT}\), \(\Delta_4\),
\(R_{\rm act}\), \(\kappa_e\), or the MeV chart. Those evaluations belong to
the arithmetic chain of the action line and to the energetic realization of
\cref{thm:reduccion-electron-two-way-rev11}; this arithmetic chain remains
separate from the combinatorial verification of the readers.

The metrological comparison subsequently applies the operators to the typed inventory,
declares unit, scheme, scale, date, and uncertainty, and performs permutations
and validation by retaining complete multisections. The physical output is
governed by the domain \(81/56/13/324\), by the two route readers, and by
the fiber operators.

\begin{statusbox}[title={Exact domain and closed physical realization}]
The atlas, the \(C_{81}\to\mathcal L_{108}\)-iteration chain of the
enriched chronological record of \(108\) iterations,
the action line,
the readers \(K_D,K_\Omega\), their electron reduction, and the fibre
operators are exact internal results. The pair
\(\mathbf K=(K_D,K_\Omega)\) is the complete output on each route and
determines the two diagonal operators of every fibre. In the electron
fibre, both readers coincide. In the noncentral fibres, the pair and its
spectrum preserve the complete refinement, while
\(\mathcal S_{\rm phys}\) binds branch, generation, fine signature, and, when
applicable, colour orbit to the persistent route. The character of that route
publishes the scalar coordinate before metrological comparison. Hence there is no
disjunction between a closed operator and a pending mass: they are two typed
readings of the same genealogy. Both readers are co-governing, and no
proximity to an external mass intervenes in the selection.
\end{statusbox}
\HMTClaim{MD-R-MASS-PHYSICAL-SELECTION-CLOSED}
